\documentclass[10pt,a4paper]{amsart}
\usepackage{amsmath,amssymb,amsthm,graphicx}
\usepackage{verbatim}
\newenvironment{reference}{\comment}{\endcomment}
%\newenvironment{reference}{}{}
\newenvironment{slogan}{\comment}{\endcomment}
\newenvironment{history}{\comment}{\endcomment}

% For commutative diagrams we use Xy-pic
\usepackage[all]{xy}

% We use 2cell for 2-commutative diagrams.
\xyoption{2cell}
% Theorem environments.
%
\theoremstyle{plain}
\newtheorem{theorem}[subsection]{Theorem}
\newtheorem{proposition}[subsection]{Proposition}
\newtheorem{lemma}[subsection]{Lemma}

\theoremstyle{definition}
\newtheorem{definition}[subsection]{Definition}
\newtheorem{example}[subsection]{Example}
\newtheorem{exercise}[subsection]{Exercise}
\newtheorem{situation}[subsection]{Situation}

\theoremstyle{remark}
\newtheorem{remark}[subsection]{Remark}
\newtheorem{remarks}[subsection]{Remarks}

\numberwithin{equation}{subsection}

% Macros
%
\def\lim{\mathop{\mathrm{lim}}\nolimits}
\def\colim{\mathop{\mathrm{colim}}\nolimits}
\def\Spec{\mathop{\mathrm{Spec}}}
\def\Hom{\mathop{\mathrm{Hom}}\nolimits}
\def\Ext{\mathop{\mathrm{Ext}}\nolimits}
\def\SheafHom{\mathop{\mathcal{H}\!\mathit{om}}\nolimits}
\def\SheafExt{\mathop{\mathcal{E}\!\mathit{xt}}\nolimits}
\def\Sch{\mathit{Sch}}
\def\Mor{\mathop{\mathrm{Mor}}\nolimits}
\def\Ob{\mathop{\mathrm{Ob}}\nolimits}
\def\Sh{\mathop{\mathit{Sh}}\nolimits}
\def\NL{\mathop{N\!L}\nolimits}
\def\CH{\mathop{\mathrm{CH}}\nolimits}
\def\proetale{{pro\text{-}\acute{e}tale}}
\def\etale{{\acute{e}tale}}
\def\QCoh{\mathit{QCoh}}
\def\Ker{\mathop{\mathrm{Ker}}}
\def\Im{\mathop{\mathrm{Im}}}
\def\Coker{\mathop{\mathrm{Coker}}}
\def\Coim{\mathop{\mathrm{Coim}}}

% Boxtimes
%
\DeclareMathSymbol{\boxtimes}{\mathbin}{AMSa}{"02}

%
% Macros for moduli stacks/spaces
%
\def\QCohstack{\mathcal{QC}\!\mathit{oh}}
\def\Cohstack{\mathcal{C}\!\mathit{oh}}
\def\Spacesstack{\mathcal{S}\!\mathit{paces}}
\def\Quotfunctor{\mathrm{Quot}}
\def\Hilbfunctor{\mathrm{Hilb}}
\def\Curvesstack{\mathcal{C}\!\mathit{urves}}
\def\Polarizedstack{\mathcal{P}\!\mathit{olarized}}
\def\Complexesstack{\mathcal{C}\!\mathit{omplexes}}
% \Pic is the operator that assigns to X its picard group, usage \Pic(X)
% \Picardstack_{X/B} denotes the Picard stack of X over B
% \Picardfunctor_{X/B} denotes the Picard functor of X over B
\def\Pic{\mathop{\mathrm{Pic}}\nolimits}
\def\Picardstack{\mathcal{P}\!\mathit{ic}}
\def\Picardfunctor{\mathrm{Pic}}
\def\Deformationcategory{\mathcal{D}\!\mathit{ef}}
\usepackage{cleveref}
\crefname{originallemma}{original lemma}{original lemmas}
\Crefname{originallemma}{Original lemma}{Original lemmas}
\makeatletter
\newlabel{dualizing-lemma-dualizing-unique}{{0A7F}{source}{}{sourcepointer-1}{}}
\newlabel{dualizing-lemma-dualizing-unique@cref}{{[originallemma][0][]0A7F}{[1][0][]source}{}{}{}}
\makeatother
\crefname{originalreference}{original reference}{original references}
\Crefname{originalreference}{Original reference}{Original references}
\makeatletter
\newlabel{lemma-affine-duality}{{0A88}{source}{}{sourcepointer-2}{}}
\newlabel{lemma-affine-duality@cref}{{[originalreference][0][]0A88}{[1][0][]source}{}{}{}}
\makeatother
\title{Uniqueness of the dualizing complex on an algebraic space up to an invertible twist}
\author{The Stacks Project Authors}
\date{stacks-project:spaces-duality.tex:local-snapshot-20261012; GNU FDL 1.2 or later}
\begin{document}
\maketitle
\noindent\textit{Editorial note.} All delivered mathematical text is the author's own, copied from the pinned Stacks Project chapter file spaces-duality.tex; no statement or proof step was rewritten, reordered or removed.\par\medskip
\noindent\textit{Editorial note.} The counted proof is the author's complete argument: the definition of the candidate invertible object as an internal hom, and the local reduction verifying invertibility and the evaluation isomorphism.\par\medskip
\noindent\textit{Editorial note.} The label lemma-dualizing-unique-spaces names the counted statement, and the proof environment immediately following it in the source is delivered verbatim as its proof.\par\medskip
\noindent\textit{Editorial note.} The counted statement is carried by the author's own lemma environment in the chapter source.\par\medskip
\noindent\textit{Editorial note.} Cross-chapter same-name check: the label lemma-dualizing-unique-spaces is not carded in any other chapter. Pre-open gate: python operations/agent-authoring/tools/precard\_gate.py --source spaces-duality.tex --label lemma-dualizing-unique-spaces --start 463 --end 483 --local-id 28175 returned RESULT=OK.\par\medskip
\noindent\textit{Editorial note.} The statement and proof are copied from the pinned chapter file; the cited inputs are kept as external pointers with their canonical Stacks tags.\par\medskip
\noindent\textit{Allowed standard prerequisites (editorial).}\begin{itemize}
\item Affine duality over an affine etale chart. Same chapter, the lemma comparing the duality functor on an affine etale chart with the algebraic duality functor. Reduces the statement to the corresponding statement over a ring.
\item Uniqueness of dualizing complexes over a ring. Dualizing Complexes, the lemma that two dualizing complexes over a ring differ by an invertible complex, together with its proof. Shows that the internal hom between the two complexes is invertible and that the evaluation map is an isomorphism, affine locally.
\end{itemize}\medskip
\ifdefined\hypertarget\hypertarget{sourcepointer-1}{}\fi
\noindent\textit{Source reference.} Original lemma 0A7F: sources/stacks-chow/tags.txt, source lines 12878--12878. Prerequisite pointer invoked at proof line(s) 480; it names a standard fact of a sibling chapter that the counted argument uses, and every step of the counted proof built on it is retained verbatim in the delivered body.\par\medskip
\ifdefined\hypertarget\hypertarget{sourcepointer-2}{}\fi
\noindent\textit{Source reference.} Original reference 0A88: sources/stacks-spaces-duality/pinned-original-source.tex, source lines 307--307. Prerequisite pointer invoked at proof line(s) 478; it names a standard fact of the same chapter that the counted argument uses, and every step of the counted proof built on it is retained verbatim in the delivered body.\par\medskip
\begin{lemma}
\label{lemma-dualizing-unique-spaces}
Let $S$ be a scheme.
Let $X$ be a locally Noetherian algebraic space over $S$.
If $K$ and $K'$ are dualizing complexes on $X$, then $K'$
is isomorphic to $K \otimes_{\mathcal{O}_X}^\mathbf{L} L$
for some invertible object $L$ of $D(\mathcal{O}_X)$.
\end{lemma}

\begin{proof}
Set
$$
L = R\SheafHom_{\mathcal{O}_X}(K, K')
$$
This is an invertible object of $D(\mathcal{O}_X)$, because affine locally
this is true. Use Lemma \ref{lemma-affine-duality} and
Dualizing Complexes, Lemma
\ref{dualizing-lemma-dualizing-unique} and its proof.
The evaluation map $L \otimes_{\mathcal{O}_X}^\mathbf{L} K \to K'$
is an isomorphism for the same reason.
\end{proof}

\end{document}
